\documentclass[10pt,a4paper]{amsart}
\usepackage[margin=19mm]{geometry}
\usepackage{amsmath,amssymb,mathtools}
\usepackage[scr=rsfs,cal=euler]{mathalfa}
\usepackage{xfrac}
\usepackage[unicode,hidelinks,bookmarks=false]{hyperref}
\newcommand{\ie}{i.e.\ }
\newcommand{\cf}{cf.\ }
\newcommand{\resp}{respectively\ }
\newcommand{\Subsection}{\subsection}
\providecommand{\nobreakdash}{\nobreak\hbox{-}}
\setlength{\emergencystretch}{2em}
\multlinegap0pt
\usepackage{bm}
\usepackage{bbm}
\usepackage{dsfont}
\usepackage{xspace}
\usepackage{cancel}
\usepackage{mathtools}
\usepackage{amsthm}
\makeatletter
\@ifundefined{lemma}{\newtheorem{lemma}{Lemma}}{}
\@ifundefined{proposition}{\newtheorem{proposition}{Proposition}}{}
\@ifundefined{theorem}{\newtheorem{theorem}[lemma]{Theorem}}{}
\makeatother
\title[The State-Dependent Riccati Equation in Nonlinear Optimal Co]{The State-Dependent Riccati Equation in Nonlinear Optimal Control: Analysis, Error Estimation and Numerical Approximation}
\author{Luca Saluzzi}
\date{Source: arXiv:2503.01587v2 (CC BY 4.0)}
\begin{document}
\maketitle

\noindent\emph{Source and licence.} The State-Dependent Riccati Equation in Nonlinear Optimal Control: Analysis, Error Estimation and Numerical Approximation. Version arXiv:2503.01587v2 (CC BY 4.0), distributed under the Creative Commons Attribution 4.0 International licence, as recorded in the arXiv licence metadata for this exact version and shown on the abstract page https://arxiv.org/abs/2503.01587v2. Full rights notice: licenses/arxiv-2503.01587-CC-BY-4.0.txt.\par\smallskip

\noindent\textit{Editorial scope.} Counted unit: the Proposition whose source label is \texttt{eq:bound} in the article (source0.tex lines 566--593, source label \texttt{eq:bound}), delivered together with the article's own complete original proof of it (source0.tex lines 595--656, one \texttt{proof} environment of 62 lines). No mathematics is written, repaired or re-derived: statement and proof are the article's own text line for line. Required context retained uncounted: eq (lines 130--135); optc (lines 160--162); control\_sdre (lines 228--231); thm:stable (lines 259--263); utilde\_sdre (lines 277--280). Every definition, display and label of those lines is retained unchanged. Omitted from this package and not redistributed: 1--129, 136--159, 163--227, 232--258, 264--276, 281--565, 594--594, 657--1346. Nothing inside a retained excerpt is omitted or altered: every retained body is delivered exactly as the article's lines are written, so no omission receipt of any kind accompanies this package. Citations: the retained excerpts carry no citation command at all, so no bibliography entry of the article is transcribed and no reference is redistributed. Reference adaptation: none was needed and none was made; every reference inside the delivered unit resolves inside it, so no unresolved reference or citation is produced. Printed slips kept literal: no printed word, symbol, hypothesis or number of the counted statement or of its proof was altered. Layout and class: the article's own class is \texttt{article} with packages including geometry, graphicx, caption, enumerate, comment, mathtools, amsmath, amssymb, algorithm, amsthm, algpseudocode, xcolor; the wrapper is the lane's pinned \texttt{amsart} editorial wrapper at 10pt on A4, which restates the theorem-like environments the retained text needs and adds the article's own macro definitions that the retained text needs (none), with extra wrapper packages (bm, bbm, dsfont, xspace, cancel, mathtools). The delivered numbering is the unit's own. Assets: no retained span contains a figure environment, a graphics-inclusion command, a PDF-page inclusion, a table or a TikZ picture; no image file is copied into this package, the package declares no asset dependency and no image file is redistributed with it. Licence: version arXiv:2503.01587v2 of this article is distributed under the Creative Commons Attribution 4.0 International licence, as recorded in the arXiv licence metadata for this exact version and shown on the abstract page https://arxiv.org/abs/2503.01587v2. A full rights notice is delivered at licenses/arxiv-2503.01587-CC-BY-4.0.txt. Characters: every retained excerpt is pure ASCII, so the delivered engine, XeLaTeX with its default Latin Modern text font, reports no missing character.
\begin{equation}\label{eq}
\left\{ \begin{array}{l}
\dot{y}(s) = f(y(s)) + B(y(s)) u(s), \quad s \in (0, +\infty),\\
y(0) = x \in \mathbb{R}^d.
\end{array} \right.
\end{equation}  

\begin{equation}\label{optc}
u^*(x) = -\frac{1}{2} R^{-1} B(x)^{\top} \nabla V(x)\,,
\end{equation}  

\begin{equation}
	u_S(y) = -R^{-1} B^{\top}(y) P(y)y\,,
	\label{control_sdre}
\end{equation}

\begin{theorem}
    Assume that the system \eqref{eq} is such that \( f(x) \) and \( \frac{\partial f(x)}{\partial x_j} \) (for \( j = 1, \dots, n \)) are continuous in \( x \) for all \( \|x\| \leq \hat{r} \), where \( \hat{r} > 0 \).
Assume further that \( A(x) \) and \( B(x) \) are continuous and that the pair $(A(x),B(x))$ is stabilizable and the pair $(A(x),Q^{1/2})$ is detectable in some nonempty neighborhood of the origin \( \mathcal{O} \subseteq B_{\hat{r}}(0) \). Then, the system with the control given by \eqref{control_sdre} is locally asymptotically stable.
\label{thm:stable}
\end{theorem}

\begin{equation}
\tilde{u}_{S}(x) = -\frac{1}{2} R^{-1} B(x)^T \nabla V_S (x),
\label{utilde_sdre}
\end{equation}

\begin{proposition}
Under the assumptions of Theorem \ref{thm:stable}, the following error bound holds in some nonempty neighborhood of the origin
    \begin{equation}
|V_S(x) - V(x)| \leq \max\Bigl\{\int_0^{\infty} |E(y^*(t))| dt, \int_0^{\infty} |E(\tilde{y}_{S}(t))| dt\Bigr\}, 
\label{eq:bound}
\end{equation}
where $y^*(t)$ denotes the state trajectory corresponding to the optimal control \eqref{optc}, and $\tilde{y}_{S}(t)$  is the trajectory associated with the control \eqref{utilde_sdre}.
%\begin{equation}
%\tilde{u}=-\frac{1}{2} R^{-1} B(x)^\top \nabla V_S(x).
%\label{eq:utilde}
%\end{equation}

Furthermore, assume that both controlled trajectories are exponentially stable and remain in a neighborhood $\Omega$ of the origin. Then there exists a constant $C>0$ such that
\begin{equation}
|V_S(x)-V(x)|
\le
\frac{C}{\lambda}\,\|E\|_{*,\Omega}\,\|x\|,
\qquad x\in\Omega,
\label{eq:error_weighted}
\end{equation}
where
\[
\|E\|_{*,\Omega}:=\sup_{z\in\Omega\setminus\{0\}} \frac{|E(z)|}{\|z\|}.
\]



\end{proposition}

\begin{proof}


Applying the DPP for the optimal value function, we obtain that for any \(T\ge 0\)
\begin{equation}\label{eq:Vopt}
V(x)=\int_0^T \ell\bigl(y^*(s),u^*(s)\bigr)ds+V\bigl(y^*(T)\bigr),
\end{equation}
where \(y^*(t)\) is the state trajectory corresponding to the optimal control \(u^*(t)\).

Now, applying the same control \(u^*(\cdot)\) into the DPP for the SDRE approximation, we obtain
\begin{equation}\label{eq:VSDRE}
V_S(x)\le \int_0^T \Bigl[\ell\bigl(y^*(s),u^*(s)\bigr)+E\bigl(y^*(s)\bigr)\Bigr]ds+V_S\bigl(y^*(T)\bigr).
\end{equation}
By subtracting equation \eqref{eq:Vopt} from equation \eqref{eq:VSDRE} and utilizing the local asymptotic stability of the trajectory, the terminal terms vanish in the limit as $T \rightarrow + \infty$, leading to
\[
V_S(x)-V(x)\le \int_0^\infty E\bigl(y^*(s)\bigr)ds.
\]
%yields
%$$
%V_S(x)-V(x)\le  \int_0^T E\bigl(y^*(s)\bigr)ds+ \Bigl[V_S\bigl(y^*(T)\bigr)-V\bigl(y^*(T)\bigr)\Bigr].
%$$
%Since the trajectory is local asymptotically stable, the terminal terms vanish in the limit as $T \rightarrow + \infty$, leading to

A similar bound for $V(x)-V_S(x)$ can be derived by considering the control $\tilde{u}_{S}$ given in \eqref{utilde_sdre}, obtaining the bound \eqref{eq:bound}.

%Now, suppose further that the closed-loop systems under \( u^*(x) \) and $\tilde{u}(x)$ are exponentially stable. Then there exist a decay rate \( \lambda > 0 \) and a constant \( C > 0 \) such that along the trajectory,  
%\[
%E(y^*(t)) \leq C \|E\|_{\infty} e^{-\lambda t},
%\]
%so that  
%\[
% \int_0^\infty E(y^*(t)) \, dt  \leq \int_0^\infty C \|E\|_{\infty} e^{-\lambda t} \, dt = \frac{C}{\lambda} \|E\|_{\infty}.
%\]
%The same bound holds for the term involving $E(\tilde{y}_{S}(t))$, completing the proof.

Now, by definition of $\|E\|_{*,\Omega}$,
\[
|E(z)|\le \|E\|_{*,\Omega}\|z\|,
\qquad z\in\Omega.
\]
Hence, using exponential stability,
\[
|E(y^*(t;x))|
\le
\|E\|_{*,\Omega}\|y^*(t;x)\|
\le
C \|E\|_{*,\Omega} e^{-\lambda t}\|x\|,
\]
and similarly for $\tilde y_S(t;x)$. Therefore,
\[
\int_0^\infty |E(y^*(t;x))|\,dt
\le
C \|E\|_{*,\Omega}\|x\|
\int_0^\infty e^{-\lambda t}\,dt
=
\frac{C}{\lambda}\|E\|_{*,\Omega}\|x\|.
\]
The same estimate holds along $\tilde y_S(t;x)$, which proves
\eqref{eq:error_weighted}.

 
\end{proof}

\end{document}
